% !TeX root = ../paper.tex

\section{Limitations and Future Work}
\label{sec:limitations}

The findings hold under a specific set of choices. This section states the most
consequential of those bounds, grouped by the kind of validity each threatens. Two of these, the unsuitability of the Maintainability
Index and the limited power of the correctness test, are load-bearing enough that the conclusions are already phrased
around them.

\subsection{Construct Validity: Does the Measurement Capture Maintainability?}

\paragraph{The Maintainability Index.} Beyond the validity critique already
discussed~\citep{heitlager2007maintainability}, the Index proved in this study
uninformative rather than misleading: it moved almost identically across
configurations, whether or not comments were stripped, because its comment term
dominates the score at this code size, so it did not usefully discriminate between
the loop's outputs. A comment-normalized or volume-controlled measure would be a more
informative outcome metric.

\paragraph{Tool-specific operationalization.} Each metric is defined here by the
tool that computed it: the Index is \texttt{radon}'s normalized 0--100 derivative
rather than Coleman et al.'s unbounded
polynomial~\citep{coleman1994maintainability,radon_docs}, and cognitive complexity
is \texttt{complexipy}'s implementation of Campbell's
specification~\citep{cognitivecomplexity,complexipy_tool}. Both tools are pinned to the
versions used, so that a replication measures with the same implementation. The
results are therefore statements about these operationalizations rather than about
the constructs in the abstract.

\paragraph{Metric-level, not human, maintainability.} Maintainability is assessed
only through static metrics; perceived readability and human effort were not
measured. The Pylint counts are weaker still, being the loop's own feedback source
(\cref{sec:supp}), and no claim here rests on them; the model is shown its complexity
totals but is never given the Maintainability Index as a target, and the
cognitive-complexity result is a measured consequence of the refactoring instruction
rather than a restatement of it.

\subsection{Internal Validity: Could the Effect Have Another Cause?}

\paragraph{Code size as an alternative explanation.} The obvious rival explanation
is that the code simply became shorter; the supplementary measures argue against
it (\cref{sec:supp}), as the loop restructures control flow rather than shrinking
programs.

\paragraph{A single feedback prompt.} The prompt wording was not varied
systematically: the only variation tested was comment handling, whose effect on the
Index was not significant ($p=.13$). The study therefore provides no evidence about
sensitivity to the rest of the phrasing---which findings are shown, in what order,
and how forcefully.

\paragraph{Two factors inside the correctness gate.} The guarded configuration
changes two things at once: it rolls back a failing refactor, and it tells the model
its change altered behavior. The zero-regression guarantee follows from the rollback
alone, but the gate's effect on the complexity results cannot be attributed to
either factor separately; separating them would need a third condition that rolls
back silently.

\paragraph{Limits of the correctness oracle.} MBPP+ inputs are generated
automatically, so passing the oracle establishes agreement with the reference
solution on the tested inputs rather than correctness in general; a refactor that
breaks behavior only on an untested input would be recorded as correct. The oracle
also treats a timeout as a failure, which is the desired behavior for a
non-terminating refactor but means a genuinely slow solution could be scored the
same way. Neither limitation is specific to this study---it is inherited from the
benchmark---but both bound how strong ``still correct'' is as a claim.

\paragraph{The iteration cap.} The cap of three was fixed in advance as a budget,
not derived from a power analysis. The trajectory analysis
(\cref{sec:trajectory}) later showed it was not binding for most tasks, but that is
a post-hoc justification and does not establish that a larger budget would change
nothing.

\subsection{External Validity: How Far Does This Generalize?}

\paragraph{Scope of the sample.} The study uses a single model
(Qwen2.5-Coder-7B-Instruct) at a single fixed seed, on 63 short benchmark tasks
(23 in the higher-complexity stratum). Short, self-contained tasks have limited
explanatory power for large software systems, so all claims are scoped strictly to
this experimental setting. Whether the cognitive-complexity effect generalizes to
other model families and sizes was not tested and is left to future work.

\paragraph{One decoding configuration and one quantization.} Generation used greedy
decoding at temperature~0 and the Q4\_K\_M quantization throughout. Both were chosen
for reproducibility, and both plausibly matter: quantization may affect output
quality, and sampling at a higher temperature would introduce the run-to-run
variation that temperature~0 deliberately removes. Neither was varied, so the
results describe one point in that space.

\paragraph{Model identity.} An Ollama tag names a model rather than fixing its
content. The digests of the weights used are recorded (\cref{sec:method}) so a
replication can check them, but they were read back after the fact, and the evidence
that one set of weights served all five runs is timestamp-based.

\subsection{Statistical Conclusion Validity}

\paragraph{Power and multiplicity.} A family of twelve tests is
reported (the three primary metrics and the correctness test, each across the
three configurations). To control the family-wise error rate, a Bonferroni
correction is applied, giving an adjusted threshold of
$\alpha^\ast = 0.05/12 = 0.00417$. The cognitive-complexity reductions and the
Maintainability-Index movements remain significant under this correction in every
configuration; only the weaker cyclomatic-complexity results in the guarded and
comment-preserving loops fall away. Bonferroni controls the family-wise error rate
at $\alpha = 0.05$ regardless of the dependence among the tests. Several outcomes
here are computed from the same tasks and the configurations share a common
baseline, so the tests are not independent; the correction may consequently be
conservative, although the degree of that conservatism is not quantified here. A
step-down (Holm) correction~\citep{holm1979} rejects the
same set. The
code-blind ablation (\cref{sec:blind}) is kept in a separate secondary family of
eight tests with its own threshold ($\alpha^\ast = 0.05/8 = 0.00625$), under
which both code-blind cognitive-complexity reductions survive; were the two
families instead pooled into a single set of twenty tests
($\alpha^\ast = 0.05/20 = 0.0025$), the guarded code-blind result ($p=.005$) and the
comment-stripping cyclomatic-complexity result ($p=.0036$) would both cease to be significant; the unguarded code-blind result ($p<.001$)
would still clear it, and the cognitive-complexity and Maintainability-Index
conclusions for the three primary configurations are unchanged. The per-stratum
analyses are treated as secondary and exploratory. Independently, with 63 tasks
and only a handful of discordant pairs the correctness (McNemar) test has low
power, so the non-significant downward trend of the unguarded loop should not be
over-interpreted.

\subsection{Reproducibility and Artifact Availability}
\label{sec:repro}

\paragraph{What is reproducible, and how exactly.} As set out in
\cref{sec:artifacts}, recomputation and re-running are different claims. All 1260
retained solutions re-measure to their recorded values, so each figure traces to a
specific artifact and the statistics can be re-derived offline. Re-running
generation is weaker: the baseline is byte-identical across configurations, but a
multi-step trajectory can diverge once any intermediate iteration differs. The
quantities above are therefore measurements of a single archived run per
configuration, not constants a fresh execution will reproduce digit for digit.

\paragraph{An incident that shaped this position.} The distinction just drawn was
learned rather than adopted on principle. Two loop runners originally wrote their
generated code to the same filename pattern, so the second silently overwrote the
first run's saved solutions; the overwrite was found by an integrity audit before
submission and the affected configurations were re-run under collision-free names.
On the fresh runs a substantially larger Maintainability-Index effect previously
recorded for the comment-stripping configuration did not reproduce; it fell from roughly $-16$ points to the $-3.0$ reported here, and the strip-versus-preserve
difference became non-significant. That figure and the ablation argument built on it
were withdrawn, and the fresh runs adopted as canonical. The instability was
specific to multi-step trajectories, which is why this thesis rests its conclusions
on effect direction and significance rather than on precise magnitudes.

\subsection{Future Work}

The bounds above suggest the following directions, roughly in order of how directly
they address a limitation identified here.

\begin{itemize}
  \item \textbf{Replicate across models.} Repeat the study on structurally messier
    code generators to test whether the cognitive-complexity effect is a property of
    the technique or of this model.
  \item \textbf{Vary the prompt and pre-register.} Measure sensitivity to which
    findings are shown and how the instruction is phrased---the factor held fixed
    here---and fix the metric family, correction scheme, and iteration budget in
    advance.
  \item \textbf{Adopt a volume-controlled maintainability measure.} A metric that
    does not conflate code size and comment density with structural quality would
    give H1 an informative outcome variable instead of one at its ceiling.
  \item \textbf{Add human raters.} Test whether the measured reductions correspond
    to code developers find easier to read and change, and decompose the gate with a
    silent-rollback condition.
  \item \textbf{Scale beyond self-contained functions.} Apply the loop to
    repository-level code, where maintainability matters most and the
    single-function metrics used here are least adequate.
\end{itemize}
